\documentclass[11pt,a4paper]{article}
\usepackage[T1]{fontenc}
\usepackage[utf8]{inputenc}
\usepackage{lmodern}
\usepackage[margin=23mm,headheight=24pt]{geometry}
\usepackage{xcolor}
\usepackage{graphicx}
\usepackage{booktabs}
\usepackage{longtable}
\usepackage{tabularx}
\usepackage{array}
\usepackage{enumitem}
\usepackage{amsmath}
\usepackage{microtype}
\usepackage{fancyhdr}
\usepackage{titlesec}
\usepackage{listings}
\usepackage{hyperref}
\usepackage{bookmark}

\definecolor{Navy}{HTML}{17324D}
\definecolor{Teal}{HTML}{087E8B}
\definecolor{Coral}{HTML}{D65A4A}
\definecolor{Pale}{HTML}{EDF4F5}
\definecolor{Ink}{HTML}{20262D}
\definecolor{SoftGray}{HTML}{F3F5F6}
\definecolor{RuleGray}{HTML}{CBD2D7}

\hypersetup{
  colorlinks=true,
  linkcolor=Navy,
  urlcolor=Teal,
  citecolor=Teal,
  pdftitle={3D-BrAIn User and Scientific Manual},
  pdfauthor={Cristina Campi},
  pdfsubject={High-density 3Brain MEA data analysis}
}

\pagestyle{fancy}
\fancyhf{}
\lhead{\textcolor{Navy}{3D-BrAIn}}
\rhead{\textcolor{Navy}{User and Scientific Manual}}
\cfoot{\thepage}
\renewcommand{\headrulewidth}{0.4pt}
\renewcommand{\headrule}{\hbox to\headwidth{\color{Teal}\leaders\hrule height \headrulewidth\hfill}}

\titleformat{\section}{\Large\bfseries\color{Navy}}{\thesection}{0.7em}{}
\titleformat{\subsection}{\large\bfseries\color{Teal}}{\thesubsection}{0.7em}{}
\titleformat{\subsubsection}{\normalsize\bfseries\color{Ink}}{\thesubsubsection}{0.7em}{}
\setlist{nosep,leftmargin=6mm}
\setlength{\parindent}{0pt}
\setlength{\parskip}{5pt}
\setlength{\emergencystretch}{2em}
\renewcommand{\arraystretch}{1.2}

\lstdefinestyle{terminal}{
  basicstyle=\ttfamily\small,
  backgroundcolor=\color{SoftGray},
  frame=single,
  rulecolor=\color{RuleGray},
  breaklines=true,
  columns=fullflexible,
  keepspaces=true,
  showstringspaces=false,
  xleftmargin=2mm,
  xrightmargin=2mm,
  aboveskip=7pt,
  belowskip=7pt
}
\lstset{style=terminal}

\newcommand{\code}[1]{\texttt{#1}}
\newcommand{\note}[1]{%
  \begin{center}
  \fcolorbox{Teal}{Pale}{\begin{minipage}{0.92\linewidth}\textbf{Method note.} #1\end{minipage}}
  \end{center}
}
\newcommand{\warningbox}[1]{%
  \begin{center}
  \fcolorbox{Coral}{SoftGray}{\begin{minipage}{0.92\linewidth}\textbf{Quality-control warning.} #1\end{minipage}}
  \end{center}
}

\begin{document}

\hypersetup{pageanchor=false}
\begin{titlepage}
  \pagecolor{Navy}
  \color{white}
  \vspace*{22mm}
  {\Huge\bfseries 3D-BrAIn\par}
  \vspace{5mm}
  {\LARGE User and Scientific Manual\par}
  \vspace{7mm}
  {\large Reproducible analysis of high-density 3Brain MEA recordings\par}
  \vfill
  \begin{tabular}{@{}ll@{}}
    Software version & 1.0.0 \\
    Manual edition & 10 September 2026 \\
    Author & Cristina Campi \\
    Repository & \href{https://github.com/cristinacampi/3D-BrAIn_codes}{github.com/cristinacampi/3D-BrAIn\_codes} \\
    License & MIT \\
  \end{tabular}
  \vspace{16mm}

  {\small Research software. Not a clinical or diagnostic device.}
\end{titlepage}
\nopagecolor
\color{Ink}
\hypersetup{pageanchor=true}
\pagenumbering{roman}

\section*{Document Purpose}

This manual is the publication-oriented reference for 3D-BrAIn. It combines a
user guide, computational methods description, output dictionary, model-reuse
guide, and reproducibility checklist. It describes the current public command
line programs and their default settings. Dataset-specific reports should add
the actual recording counts, biological units, DIV values, conditions, retained
channels, retained spikes, and measured results.

\warningbox{Defaults are conveniences, not acquisition metadata. Confirm the
sampling rate, duration, well, grid convention, reference channel, voltage
scale, and labels for every dataset.}

\tableofcontents
\clearpage
\pagenumbering{arabic}

\section{Software Scope and Architecture}

3D-BrAIn is a Python toolkit for high-density microelectrode-array data from
3Brain systems. It supports raw file access, signal extraction, event detection,
spike and burst metrics, waveform stratification, recording-level supervised
learning, functional connectivity, and reusable research models.

The project deliberately separates orchestration from scientific logic:

\begin{longtable}{>{\ttfamily}p{0.22\linewidth}p{0.70\linewidth}}
\toprule
Location & Responsibility \\
\midrule
main/ & Configurable analysis programs containing orchestration and editable parameters. \\
src/brain\_3d/ & Scientific functions, preprocessing, models, figures, metrics, and output generation. \\
tests/ & Deterministic tests for API conventions, preprocessing, connectivity, Mamba, and configurable scripts. \\
docs/ & GitHub guides, Sphinx API documentation, methods, and printable source. \\
tools/ & Report repair/building and model-portal integration utilities. \\
\bottomrule
\end{longtable}

Public scientific modules include \code{BrwFunctions}, \code{BxrFunctions},
\code{SpikeSorting}, \code{Stratification},
\code{Classification}, \code{Regression}, \code{Connectivity}, and
\code{Mamba}. Public functions and classes use the project's PascalCase
convention throughout the scientific modules.

\section{Installation}

Python 3.8 or newer is required. An isolated environment is recommended.

\begin{lstlisting}
git clone https://github.com/cristinacampi/3D-BrAIn_codes.git
cd 3D-BrAIn_codes
python3 -m venv .venv
source .venv/bin/activate
python3 -m pip install --upgrade pip
python3 -m pip install -e .
\end{lstlisting}

Install development and documentation dependencies with:

\begin{lstlisting}
python3 -m pip install -e ".[dev]"
python3 -m pytest -q
\end{lstlisting}

The repository also provides \code{environment.yml}, \code{Dockerfile}, and
\code{docker-compose.yml}. GPU acceleration depends on a PyTorch build that is
compatible with the local CUDA environment. The official \code{mamba-ssm}
backend is optional and must be installed separately on a compatible platform;
the default reference implementation uses standard PyTorch operations.

\section{Input Data Contracts}

\subsection{Raw BRW data}

Trace extraction expects a directory containing one or more BRW recordings.
The default well is \code{Well\_A1}. Sampling rate, frame count, and channel
metadata are read from the raw recording and stored with extracted HDF5 traces.

\subsection{Processed spike arrays}

Each immediate subdirectory of a spike root represents one recording. Matched
files are expected for each channel \(c\):

\begin{longtable}{>{\ttfamily}p{0.28\linewidth}p{0.22\linewidth}p{0.42\linewidth}}
\toprule
File & Shape & Meaning \\
\midrule
spikes\_ch\(c\).npy & events \(\times\) samples & Waveform around each detected event. \\
frames\_ch\(c\).npy & events & Detection frame for each event. \\
label\_ch\(c\).npy & events & Detector polarity: 0 negative, 1 positive. \\
\bottomrule
\end{longtable}

When matched arrays differ in length, analysis uses their shortest common event
count. This protects indexing but signals an upstream data-integrity problem
that should be corrected or reported. Mamba recording sequences additionally
require equal waveform length across recordings.

Channel 0 is excluded as the reference by negative-waveform workflows. On a
64 by 64 grid, 4095 channels remain before activity and waveform filters. The
default zero-based mapping is
\[
  r = \lfloor c/64 \rfloor, \qquad q = c \bmod 64.
\]
Verify this convention against acquisition metadata before interpreting sensor
maps.

\subsection{Labels table}

Supervised analyses expect one CSV row per recording. Recording matching is
case-insensitive after stripping whitespace and a terminal \code{.brw}
extension.

\begin{longtable}{>{\ttfamily}p{0.26\linewidth}p{0.64\linewidth}}
\toprule
Column & Purpose \\
\midrule
filename & Recording file or recording-folder name. \\
biological\_unit\_id & Organoid, culture, subject, or other unit that must never cross validation folds. \\
chip\_ID & Chip identifier retained in analysis metadata. \\
DIV & Numeric days in vitro; continuous target for regression. \\
Days\_in\_BF & Experimental metadata retained by current supervised pipelines. \\
culture\_phase\_code & Numeric class code for the built-in culture-phase task. \\
culture\_phase & Human-readable culture-phase label. \\
div\_code, div\_label & Optional discrete DIV class and readable label. \\
experiment\_id & Required only when the labels table is filtered by experiment. \\
\bottomrule
\end{longtable}

Repeated measurements from the same organoid must carry the same
\code{biological\_unit\_id}. Otherwise, information from one organoid can leak
between training and test folds.

\section{End-to-End Workflow}

\subsection{Trace extraction}

\begin{lstlisting}
python3 main/main_extract_traces_h5.py
\end{lstlisting}

Edit the parameter block before running. The program exports chunked HDF5
traces. Restrict the channel range with the configured start and end values.
Existing complete extractions are skipped unless recomputation is enabled.

\subsection{Spike detection}

\begin{lstlisting}
python3 main/main_detect_spikes_from_h5.py
\end{lstlisting}

The waveform length must be odd so the event can be centred. Default filter and
detection settings must be reported and should be validated against voltage
units, noise, sampling rate, and expected extracellular waveform duration.

\subsection{Metrics}

\begin{lstlisting}
python3 main/main_compute_spike_burst_metrics.py
\end{lstlisting}

Set the activity threshold to 0.25 Hz in the parameter block to align it with
the negative-waveform analyses. The metric output root is also configured in
the same block.

\subsection{Analysis branches}

After spike extraction, the same processed folder can feed four complementary
branches:

\begin{enumerate}
  \item unsupervised waveform stratification;
  \item recording-level classification;
  \item recording-level continuous DIV regression;
  \item functional connectivity from spike times.
\end{enumerate}

All four negative-spike branches use the shared event selection described next.

\section{Shared Plausible-Negative Filter}

For waveform matrix \(X\), event \(i\) is retained when the following active
criteria are satisfied:

\begin{enumerate}
  \item detector label \(y_i=0\), when finite labels are available;
  \item \(\lvert\min X_i\rvert \geq 1.10\lvert\max X_i\rvert\);
  \item the trough index lies within 10 samples of the waveform centre.
\end{enumerate}

The detector-label check is skipped if no finite labels are available, but the
waveform checks remain active. These rules reject positive-dominant and poorly
centred events. They are operational plausibility criteria, not proof that an
event originates from a neuron or a universal definition of extracellular spike
polarity.

After event filtering, channel activity is
\[
  f_c = \frac{N_c}{T},
\]
where \(N_c\) is the retained event count and \(T\) is recording duration in
seconds. A channel is active when \(f_c > 0.25\) Hz by default. The strict
inequality means a channel exactly at 0.25 Hz is excluded.

\note{Every publication should report input events, detector-negative events,
waveform-plausible events, active channels, and events finally sampled. These
counts distinguish biological sparsity from filtering or file-coverage effects.}

\section{Spike and Burst Metrics}

Per-channel and well-level spike outputs include event count and rate, mean ISI,
ISI variability, peak-to-peak amplitude, and active-electrode summaries.

A burst is a maximal sequence with at least \(N_{\min}=5\) spikes and
consecutive inter-spike intervals no greater than
\(\mathrm{ISI}_{\max}=0.1\) s. After the initial five-spike condition is met,
the sequence extends while the interval condition remains true. Reported burst
quantities include:

\begin{itemize}
  \item count and bursts per minute;
  \item mean duration and spikes per burst;
  \item mean within-burst ISI;
  \item inter-burst interval (IBI) and variability;
  \item percentage of spikes assigned to bursts.
\end{itemize}

The same algorithm is also applied to the pooled active-electrode event sequence
to produce well-level burst summaries. This is distinct from a network burst.
The API function \code{NetworkBurstMetric} in \code{BrwFunctions} identifies windows in
which bursts occur on a specified fraction of active electrodes, but the public
metrics main does not currently call it. Network-burst results must not be
reported unless that function is run explicitly with documented parameters.

\section{Waveform Stratification}

\subsection{Scope and sampling}

The pipeline produces independent clustering for each recording and a
transversal clustering pooled across recordings belonging to one experiment.
An experiment token of the form \code{NN-NN} is inferred from each spike-root
path. Roots without such a token are grouped as \code{unknown\_experiment}.

At most 15,000 filtered waveforms are selected per recording with seeded random
sampling. The transversal pool is capped at 150,000 events per experiment.
Sampling equalises computational load but does not force equal contribution by
recording. Contribution counts and proportions must therefore be inspected.

\subsection{PCA and k-means}

Waveforms are whitened and passed to \code{Stratification.Clustering} with
\code{Algo=PCA\&KM}, Euclidean distance, and Minkowski \(p=2\). PCA retains the
smallest component count \(m\) satisfying
\[
  \frac{\sum_{j=1}^{m}\lambda_j}{\sum_{j=1}^{d}\lambda_j} \geq 0.90,
\]
where \(\lambda_j\) is the variance explained by component \(j\). K-means is
evaluated for candidate counts 2, 3, 4, 5, and 6 by default.

The saved PC1--PC2 scatterplot is only a two-dimensional display. Clustering
uses the retained component space, which can contain more than two components.

\subsection{Interpretation}

Centroid plots show the cluster mean and a shaded pointwise interval of one
standard deviation. The channel count map shows where sampled events originate.
For every cluster \(k\), the local map reports
\[
  100\frac{N_{c,k}}{N_c},
\]
the percentage of analysed waveforms on channel \(c\) assigned to cluster
\(k\). The transversal recording profile analogously reports cluster
proportions within each recording.

Cluster IDs are nominal, run-specific labels. They do not imply order or direct
correspondence across separately fitted clusterings. Biological interpretation
requires waveform quality control, stability analysis, spatial context, and
association with independent experimental variables.

\section{Recording-Level Classification}

\subsection{Statistical unit and features}

One recording is one supervised-learning sample. Individual spikes are not
independent rows in the classifier. Each recording can be represented by:

\begin{itemize}
  \item activity features: counts through filters, total retained events,
  duration, firing rate, numbers of active/non-empty channels, channel count
  summaries, and channel entropy;
  \item scalar waveform summaries: mean, standard deviation, minimum, maximum,
  peak-to-peak range, and mean squared energy;
  \item MiniRocket features learned from waveform sequences;
  \item a five-channel aggregate waveform sequence for Mamba.
\end{itemize}

MiniRocket is fitted on at most 12,000 waveforms from the training fold. With
10,000 kernels, each waveform produces 10,000 responses. Responses are
aggregated within a recording by feature-wise mean and standard deviation,
giving 20,000 recording variables before any activity or scalar waveform
features are concatenated.

For Mamba, each recording becomes a sequence with five pointwise profiles:
waveform mean, standard deviation, median, 25th percentile, and 75th percentile.
This preserves the waveform time axis while controlling the number of supervised
samples.

\subsection{Algorithms}

The classical comparison includes:

\begin{itemize}
  \item most-frequent dummy baseline;
  \item ridge classifier with PCA;
  \item class-balanced logistic regression with PCA;
  \item class-balanced linear and radial-basis SVM with PCA;
  \item class-balanced random forest;
  \item class-balanced Extra Trees.
\end{itemize}

PCA pipelines standardise variables and retain 90\% explained variance. Forest
models use 300 trees by default. The optional Mamba model is a selective
state-space sequence model trained for 30 epochs by default.

\subsection{Cross-validation and leakage control}

The default validation uses four-fold \code{StratifiedGroupKFold}, repeated
three times with repeat-specific seeds. Class codes provide stratification and
\code{biological\_unit\_id} defines groups. No biological unit can occur in
both training and test sets within a fold. MiniRocket, scaling, PCA, and the
estimator are fitted using training-fold data only.

For \(K\) classes, balanced accuracy is macro-average recall:
\[
  \mathrm{BA}=\frac{1}{K}\sum_{k=1}^{K}
  \frac{\mathrm{TP}_k}{\mathrm{TP}_k+\mathrm{FN}_k}.
\]
It gives each class equal weight irrespective of prevalence. Accuracy,
per-class recall and precision, F1, support, confusion matrices, and train-test
differences should be interpreted alongside it.

\warningbox{Files and tables may call held-out folds ``test``. These are
cross-validation test folds, not an independent external test cohort. External
validation must use a new labelled dataset that did not affect model choice or
parameter selection.}

\section{Continuous DIV Regression}

DIV regression uses the same recording unit, event filter, activity features,
waveform summaries, and MiniRocket aggregation as classification. It predicts
numeric DIV rather than a categorical DIV label.

Models are a mean dummy baseline, ridge, Elastic Net, linear and radial-basis
support-vector regression, random forest, and Extra Trees. PCA models retain
90\% explained variance. DIV is discretised into up to four quantile bins only
for balanced grouped fold construction; predictions and scores use the original
continuous values. If stratified grouped splitting is infeasible, the pipeline
uses \code{GroupKFold}.

Reported metrics include MAE, median absolute error, RMSE, \(R^2\), Pearson and
Spearman association, and the fraction of predictions within 5, 10, and 15 days.
Absolute errors must be interpreted against the observed DIV range and the mean
baseline. Correlation alone does not establish accurate calibration.

\section{Functional Connectivity}

\subsection{Spike-count correlation}

Each active channel's retained event train is binned into windows of width
\(\Delta_b=50\) ms by default. Pairwise Pearson correlation is calculated
between the channel count vectors. A stronger absolute coefficient means more
similar or opposing binned count variation, depending on sign. The result is
sensitive to bin width, firing rate, shared drive, non-stationarity, and
recording duration.

\subsection{Spike-time tiling coefficient}

STTC uses a coincidence half-window \(\Delta_t=20\) ms by default. For spike
trains A and B it is computed from the proportions of spikes near the other
train, \(P_A,P_B\), and proportions of recording time covered by the other
train's coincidence tiles, \(T_A,T_B\):
\[
 \mathrm{STTC}=\frac{1}{2}\left(
 \frac{P_A-T_B}{1-P_AT_B}+
 \frac{P_B-T_A}{1-P_BT_A}\right).
\]
Positive values indicate temporal association, values near zero indicate little
association, and negative values indicate avoidance relative to tiled time.
STTC was designed to reduce firing-rate dependence compared with simpler
coincidence indexes \cite{cutts2014}.

\subsection{Edges, nodes, and sensor maps}

By default an edge passes when \(\lvert w_{ij}\rvert\geq0.2\). The
\code{--signed-edges} option instead thresholds signed values. The network map
then displays the strongest 5\% of threshold-passing edges, with at most 120
edges. Unless signed selection is active, strongest means greatest
\(\lvert w_{ij}\rvert\). Edge colour retains sign.

Global node strength is
\[
  s_i=\sum_{j\ne i}\lvert w_{ij}\rvert
\]
on the full matrix. Selected-map node strength is computed from displayed edges
only. Thus a globally strong node can appear unconnected when none of its edges
survives the threshold, percentile, or edge cap.

\subsection{Sensitivity analysis}

The correlation bin, STTC window, both thresholds, sign rule, top-edge
percentage, and edge cap are public parameters. A defensible sensitivity study
should define a grid before examining the final result, for example:

\begin{longtable}{p{0.36\linewidth}p{0.50\linewidth}}
\toprule
Parameter & Example values \\
\midrule
Correlation bin width & 20, 50, 100, and 200 ms \\
STTC coincidence window & 5, 10, 20, and 50 ms \\
Absolute edge threshold & 0.1, 0.2, 0.3, and 0.4 \\
Displayed passing edges & 1\%, 5\%, 10\%, and all \\
\bottomrule
\end{longtable}

For each setting compare active-channel count, passing-edge density, positive
and negative edge fractions, node-strength rank stability, overlap of strongest
edges, and spatial concentration. Parameter exploration must be identified as
exploratory if it influences the reported setting.

\section{Public Command Reference}

\begin{longtable}{>{\ttfamily\scriptsize\raggedright\arraybackslash}p{0.54\linewidth}p{0.38\linewidth}}
\toprule
Program & Purpose \\
\midrule
main\_extract\_traces\_h5.py & Export BRW recordings to chunked HDF5 traces. \\
main\_export\_h5\_to\_channel\_npy.py & Optional one-array-per-channel trace export. \\
main\_detect\_spikes\_from\_h5.py & Detect positive and negative events and save waveforms. \\
main\_compute\_spike\_burst\_metrics.py & Compute channel and well spike/burst tables. \\
main\_prepare\_labeled\_metrics.py & Join labels with one or more metric roots. \\
main\_run\_negative\_spike\_stratification\_analysis.py & Per-recording and transversal PCA/k-means clustering. \\
main\_run\_negative\_spike\_connectivity\_analysis.py & Correlation, STTC, edge, strength, and map analysis. \\
main\_run\_negative\_spike\_recording\_classification.py & Grouped classification and reusable final models. \\
main\_run\_negative\_spike\_div\_regression\_analysis.py & Grouped continuous-DIV regression. \\
main\_apply\_saved\_negative\_spike\_classifier.py & Extract features and apply a saved classifier to new data. \\
main\_run\_mamba\_algorithm\_comparison.py & Compare Mamba and classical classifiers on NPZ sequences. \\
\bottomrule
\end{longtable}

Every program has an editable parameter block near its top. Configuration and
processing-order guidance is maintained in \code{docs/public\_cli\_usage.md}.

\section{Output Dictionary}

\subsection{Metrics}

Each recording's \code{spike\_burst/} directory includes channel and well CSV
tables plus a complete pickle object. The metric root also contains
\code{all\_recordings\_spike\_burst\_summary.csv}.

\subsection{Stratification}

Per-recording folders include \code{run\_info.csv}, spike assignments, centroid
means and standard deviations, cluster-by-channel tables, PCA and centroid
figures, event-count maps, and local cluster-percentage maps. The transversal
folder adds recording-by-cluster counts and normalised profiles.

\subsection{Classification and regression}

Supervised analyses save recording metadata, fold metrics, held-out predictions,
ranked summaries, diagnostic figures, and reusable artifacts. The
\code{final\_models/} directory includes one bundle per selected algorithm,
JSON model metadata, \code{model\_manifest.csv}, and
\code{training\_recordings.csv}.

\subsection{Connectivity}

Each recording includes correlation and STTC matrices, threshold-passing edge
lists, full node-strength tables, firing rates, filter counts, binned counts,
parameters, and network maps. The maps folder stores selected edge lists and
selected-map strength separately from global strength.

Generated outputs are excluded from Git. A publication archive should keep
tables, figures, model metadata, exact commands, logs, and checksums together.

\section{Saved Models and New Data}

After cross-validation, selected models are refitted on all eligible labelled
recordings. Their metadata records training recordings and biological units,
class mapping or target, feature set, preprocessing, estimator parameters,
random seed, and train/test cross-validation metrics.

Configure the saved model, spike root, and output folder, then run:

\begin{lstlisting}
python3 main/main_apply_saved_negative_spike_classifier.py
\end{lstlisting}

The application command repeats waveform filtering, activity selection, and
feature extraction using the model bundle. A precomputed feature CSV is not
required. New data must match training waveform length, detector labels,
sampling assumptions, voltage scale, channel convention, and experimental
domain. A successful prediction is not external validation.

Joblib, pickle, and PyTorch files may execute code when loaded. Use trusted
artifacts only and retain checksums after transfer.

\section{Reproducibility and Publication Checklist}

\subsection{Dataset and acquisition}

\begin{itemize}
  \item Record experiment, organoid/culture, chip, well, condition, DIV, date,
  hardware, sampling rate, duration, and voltage units.
  \item Preserve raw-recording and labels-file checksums.
  \item State recording and biological-unit inclusion/exclusion rules.
  \item Keep experimental campaigns separate unless pooling is pre-specified.
\end{itemize}

\subsection{Computation}

\begin{itemize}
  \item Save \code{git rev-parse HEAD}, Python version, and package environment.
  \item Archive the exact command, standard output/error log, and random seed.
  \item Report event counts at every filter stage and all sampling caps.
  \item Verify expected, present, non-empty, and active channel counts.
  \item Fit learned feature transforms only within training folds.
  \item Compare supervised results with dummy baselines and report fold spread.
  \item Distinguish grouped cross-validation from external validation.
  \item Report and test connectivity windows, thresholds, and map selection.
\end{itemize}

\subsection{Figures and release}

\begin{itemize}
  \item Inspect legends, colour bars, labels, clipping, and resolution.
  \item Keep cluster colours and IDs explicitly run-specific.
  \item Update software version, changelog, citation metadata, and manual date.
  \item Run tests and build Sphinx documentation with warnings as errors.
  \item Attach this PDF and checksums to the software release.
\end{itemize}

\section{Troubleshooting}

\begin{longtable}{p{0.33\linewidth}p{0.59\linewidth}}
\toprule
Symptom & Checks and resolution \\
\midrule
No spike metric files found & Run spike detection and the metrics main first; confirm the spike root and recording subdirectories. \\
Only a few active channels & Confirm duration and sampling rate, inspect filter-stage counts, and remember that activity requires a rate strictly greater than the threshold after polarity filtering. \\
4096-channel expectation mismatch & Channel 0 is the reference, leaving 4095 candidates; missing spike files and empty/inactive channels are reported separately. \\
Inconsistent supervised sample counts & One feature row and one target are required per recording. Rebuild recording-level features rather than passing individual spike labels to the classifier. \\
Poor multiclass DIV classification & Prefer continuous DIV regression when maturity is ordinal/continuous and class boundaries are artificial. \\
Cluster separation is only polarity & Use the plausible-negative pipeline and inspect the filter summary and centroid variability. \\
Strong node without displayed edge & Global node strength uses the complete matrix; maps show only thresholded and top-ranked edges. \\
Existing outputs are reused & Use the relevant force option only after recording why recomputation is needed. Use map-only/reuse options when only visual outputs changed. \\
Missing Python package & Activate the environment used for installation and run \code{python -m pip install -e .}; avoid mixing system Python and environment Python. \\
\bottomrule
\end{longtable}

\section{Reporting Template}

The following compact paragraph can be adapted for a manuscript after replacing
bracketed values:

\begin{quote}\small
MEA recordings (\(n=[\mathrm{recordings}]\),
\(n=[\mathrm{biological\ units}]\)) were processed with 3D-BrAIn version
[version], commit [hash]. Signals were acquired at [sampling rate] Hz for
[duration] s. Detected events were restricted to channel-wise plausible
negative waveforms using detector label 0, a negative-to-positive extremum
ratio of at least 1.10, and a trough within 10 samples of the extraction centre.
Reference channel 0 was excluded and channels with retained firing rate no
greater than 0.25 Hz were removed. [Analysis] used [feature/method details].
Supervised models were evaluated by [folds]-fold grouped cross-validation with
[repeats] repeats, keeping all recordings from one biological unit within one
partition. Connectivity was evaluated using [bin] ms spike-count correlation
and STTC with a [window] ms coincidence window; edges were defined at
[threshold] and network maps displayed [percentage] of threshold-passing edges.
\end{quote}

Add dataset-specific counts, class or DIV distributions, model-selection rules,
effect sizes, uncertainty, sensitivity analysis, and external-validation status.

\section{References}

\begin{thebibliography}{9}

\bibitem{cutts2014}
C. S. Cutts and S. J. Eglen.
\newblock Detecting pairwise correlations in spike trains: an objective
comparison of methods and application to the study of retinal waves.
\newblock \emph{Journal of Neuroscience}, 34(43):14288--14303, 2014.
\newblock \url{https://doi.org/10.1523/JNEUROSCI.2767-14.2014}.

\bibitem{minirocket}
A. Dempster, D. F. Schmidt, and G. I. Webb.
\newblock MiniRocket: A very fast (almost) deterministic transform for time
series classification.
\newblock In \emph{Proceedings of the 27th ACM SIGKDD Conference}, 2021.
\newblock \url{https://arxiv.org/abs/2012.08791}.

\bibitem{mamba}
A. Gu and T. Dao.
\newblock Mamba: Linear-time sequence modeling with selective state spaces.
\newblock 2023. \url{https://arxiv.org/abs/2312.00752}.

\bibitem{sklearn}
F. Pedregosa et al.
\newblock Scikit-learn: Machine learning in Python.
\newblock \emph{Journal of Machine Learning Research}, 12:2825--2830, 2011.

\end{thebibliography}

\section*{Software Citation}

Use the citation shown by GitHub from the repository's \code{CITATION.cff} and
cite the exact archived release used for analysis. Cite MiniRocket, STTC, Mamba,
and other methods only when they are part of the reported analysis.

\vfill
\begin{center}
\small\textcolor{Navy}{3D-BrAIn 1.0.0 \textbullet{} MIT License \textbullet{} Research use}
\end{center}

\end{document}
